% ===== 原文 pp.93--98 =====
% [原文 p.93]

其中 $q$ 表示在 $G^{(p)}$ p-形式对称性下的电荷，$t$ 表示算符 $\cO$ 所处的、关于 $G^{(p)}$ p-形式对称性的扭曲扇区。在规范化后的理论 $\cT/G^{(p)}$ 中，算符 $\cO$ 具有性质
\begin{equation}
  (q,t)=(g,\hat{g})\in G^{(p)}\times\hat{G}^{(p)}\,, \tag{4.185}
\end{equation}
其中现在的 $q$ 是在 $G^{(d-p-2)}=\hat{G}^{(p)}$（$d-p-2$）-形式对称性下的电荷，而 $t$ 是关于 $G^{(d-p-2)}$ 的扭曲扇区。

因此，规范化把 $q$ 与 $t$ 的取值互换。

% [Tutor Note: p.93 (4.185) 是 p.92 (4.184) 的规范化版本：(4.184) 中未规范化理论的 (q,t)=(\hat g,g)\in\hat G^{(p)}\times G^{(p)}，对 G^{(p)} 规范化后两个因子按对角方式互换，故 (4.185) 写成 (g,\hat g)\in G^{(p)}\times\hat G^{(p)}。两式能同时成立的隐含前提是 d=2p+2（见 p.92 的 (4.183)）：此时 G^{(d-p-2)}=\hat G^{(p)} 的电荷与 G^{(p)} 的扭曲扇区才有匹配的维数。核对方法：取 p=1、d=4 的 Z2 1-形式对称性，其对偶仍为 1-形式对称性。]

\para{SPT 相的叠加。}

在上面的讨论中，我们对扭曲扇区算符电荷的定义比较随意。一个恰当的定义要求选定一组算符，它们位于构成该 p-形式对称性的各拓扑算符的交点处。回忆一下，在讨论 SPT 相时曾遇到过这类选择。事实上，把 $G^{(p)}$ 的 SPT 相与 $\cT$ 叠加，会使这些选择相互置换，从而改变 $\cT$ 与 $G^{(p)}$ 对称性的耦合。

让我们在一个例子中看看这一点如何起作用。

\begin{examplebox}{例 4.15：从 $\SU(2)$ 到 $\SO(3)_\pm$}
考虑 4d 纯 $\SU(2)$ Yang--Mills 理论。该理论中的扭曲扇区线算符具有形式
\begin{equation}
  \cD_{j,\phi}:=\cW_j\cH_\phi\,,\qquad j\in\Z/2\,,\quad \phi:\U(1)\to\SO(3)\,, \tag{4.186}
\end{equation}
其中余特征标 $\phi$ 必须不能提升为 $\SU(2)$ 的余特征标。这些线算符的电荷有两种选择：
\begin{enumerate}[leftmargin=1.4em,itemsep=1pt,topsep=1pt]
\item $\cD_{j,\phi}$ 在 $j\notin\Z$ 时带电、在 $j\in\Z$ 时不带电；
\item $\cD_{j,\phi}$ 在 $j\in\Z$ 时带电、在 $j\notin\Z$ 时不带电。
\end{enumerate}
这两种选择通过叠加 SPT 相 (4.161) 而相互交换。

用第一种选择给出的耦合对 $\Z_2$ 1-形式对称性作规范化，得到的理论中下列线
\begin{equation}
  \cD_{j,\phi}\,,\qquad j\in\Z\,, \tag{4.187}
\end{equation}
处于对偶 1-形式对称性的未扭曲扇区。由 3.4.1 节的讨论可知，这就是 $\SO(3)_+$ 理论。

另一方面，用第二种选择给出的耦合对 $\Z_2$ 1-形式对称性作规范化，得到的理论中下列线
\begin{equation}
  \cD_{j,\phi}\,,\qquad j\notin\Z\,, \tag{4.188}
\end{equation}
处于对偶 1-形式对称性的未扭曲扇区。由 3.4.1 节的讨论可知，这就是 $\SO(3)_-$ 理论。
\end{examplebox}

% [Tutor Note: p.93 (4.186) 中 $j\in\Z/2$ 表示“半整数”（$j=n/2$，$n\in\Z$），而不是二元素群 $\Z_2$；因此下文中 $j\notin\Z$ 即 $j$ 为半整数、$j\in\Z$ 即 $j$ 为整数（(4.223) 等处同此约定）。$\cD_{j,\phi}=\cW_j\cH_\phi$ 是 Wilson 线 $\cW_j$（$j$ 标记 $\SU(2)$ 的权/自旋）与 't Hooft 线 $\cH_\phi$（由余特征标 $\phi:\U(1)\to\SO(3)$ 标记）的乘积。可核对：$\SU(2)$ 的真 Wilson 线要求 $2j\in\Z$，故 $j\in\Z/2$；本页两种电荷选择的分界正是 $j\in\Z$ 与 $j\notin\Z$。]

% [原文 p.94]

\subsubsection{规范化交换混合 't Hooft 反常与群扩张}

考虑对 p-形式对称群 $G^{(p)}$ 的一个真子群
\begin{equation}
  H^{(p)}\subseteq G^{(p)} \tag{4.189}
\end{equation}
作规范化。规范化后的理论不仅包含对偶的
\begin{equation}
  H^{(d-p-2)}=\hat{H}^{(p)} \tag{4.190}
\end{equation}
（$d-p-2$）-形式对称性，还包含剩余的
\begin{equation}
  K^{(p)}:=G^{(p)}/H^{(p)} \tag{4.191}
\end{equation}
p-形式对称性。取决于 $H^{(p)}$ 如何嵌入 $G^{(p)}$ 之中，规范化后理论的这两个对称性会带有混合 't Hooft 反常。

为了看清这一点，注意 $G^{(p)}$ p-形式对称性的背景场 $B^G_{p+1}$ 分解为
\begin{equation}
  B^G_{p+1}=i(B^H_{p+1})+\tilde{B}^K_{p+1}\,, \tag{4.192}
\end{equation}
其中 $B^H_{p+1}$ 与 $B^K_{p+1}$ 分别是未规范化理论中 $H^{(p)}$ 与 $K^{(p)}$ 的背景场，$i$ 是单射同态
\begin{equation}
  i:\,H^{(p)}\to G^{(p)}\,, \tag{4.193}
\end{equation}
而 $\tilde{B}^K_{p+1}$ 是 $K^{(p)}$ 取值的场 $B^K_{p+1}$ 到 $G^{(p)}$ 的一个提升。由 (4.107) 附近的类似讨论可知，上述关系蕴涵
\begin{equation}
  \delta B^H_{p+1}=-\mathrm{Bock}(B^K_{p+1})\,, \tag{4.194}
\end{equation}
即一般而言 $B^H_{p+1}$ 不是闭的。

对 $H^{(p)}$ 规范化之后，关系 (4.194) 变成规范场 $b^H_{p+1}$ 的
\begin{equation}
  \delta b^H_{p+1}=-\mathrm{Bock}(B^K_{p+1})\,, \tag{4.195}
\end{equation}
耦合
\begin{equation}
  2\pi\int B_{d-p-1}\cup_\eta b^H_{p+1}\,, \tag{4.196}
\end{equation}
的存在（其中 $B_{d-p-1}$ 是 $H^{(d-p-2)}$ 的背景场）意味着规范化后的理论带有一个混合 't Hooft 反常，其相应的反常理论为

% [原文 p.95]
\begin{equation}
  \cI_{d+1}=\exp\left(-2\pi\ii\,s(d-p)\int B_{d-p-1}\cup_\eta \mathrm{Bock}(B^K_{p+1})\right),\qquad s(d-p)=(-1)^{d-p}\,. \tag{4.197}
\end{equation}

% [Tutor Note: p.94 (4.194) 的 Bock 是 Bockstein 同态，来自短正合列 (4.198)：$\mathrm{Bock}:H^{p+1}(-;K^{(p)})\to H^{p+2}(-;H^{(p)})$。推导：(4.192) 中 $B^G$ 是闭的，而 $i(B^H)$ 的闭性由 $\delta B^H$ 决定，$\delta\tilde B^K$ 不能单独为零，故 $\delta B^H=-\mathrm{Bock}(B^K)$。可核对：若 (4.198) 分裂（$K^{(p)}$ 是 $G^{(p)}$ 的正规子群且 $G^{(p)}\cong H^{(p)}\times K^{(p)}$），取 $\tilde B^K$ 为到 $K^{(p)}$ 子群因子的提升即得 $\delta B^K=0$，从而 (4.194) 平凡、(4.197) 的反常消失——这正是 p.95 紧接 (4.197) 后的论断。]

% [Tutor Note: p.94 (4.196)--(4.197) 的 $\cup_\eta$ 不是普通杯积，而是原文前文定义的“$\eta$-杯积”：$\eta:G\times H\to K$ 为双同态，$G$ 取值的 $q$-上链与 $H$ 取值的 $r$-上链经 $\eta$ 给出 $K$ 取值的 $(q+r)$-上链。这里 $B_{d-p-1}$ 取值于 $H^{(d-p-2)}=\hat H^{(p)}$、$\mathrm{Bock}(B^K_{p+1})$ 取值于 $H^{(p)}$，$\eta$ 就是这两个群之间的配对（取值于 $\R/\Z$）。核对方法：数形式次数，$(d-p-1)+(p+2)=d+1$，恰为反常理论的维数。]

注意若短正合列
\begin{equation}
  0\to H^{(p)}\to G^{(p)}\to K^{(p)}\to 0\,, \tag{4.198}
\end{equation}
分裂，则该反常消失，正如例 (4.8) 之前所解释的那样。于是我们得到如下陈述。

\begin{statementbox}{陈述 4.6：规范化把群扩张转成混合反常}
对一个位于 p-形式对称群扩张 (4.198) 之中的 p-形式对称群 $H^{(p)}$ 作规范化，会导致形式为 (4.197) 的混合 't Hooft 反常。
\end{statementbox}

上面陈述的逆命题也成立，这可以通过“对偶对称性的规范化等价于解除规范化”这一事实看出。

\begin{statementbox}{陈述 4.7：规范化把混合反常转成群扩张}
对一个带有与 p-形式对称性 $K^{(p)}$ 的形式为 (4.197) 的混合 't Hooft 反常的（$d-p-2$）-形式对称群作规范化，会得到一个对偶的 p-形式对称性 $H^{(p)}$；它与剩余的 $K^{(p)}$ 对称性结合成一个扩张后的 $G^{(p)}$ p-形式对称性，使得这三个 p-形式对称性通过 (4.198) 相联系。
\end{statementbox}

\begin{examplebox}{例 4.16：$\Spin(6)$ 与 $\SO(6)$ 纯 Yang--Mills}
纯 $\SO(6)$ Yang--Mills 理论中的混合 't Hooft 反常 (4.136) 可以看作陈述 (4.6) 的一个例子，因为 $\SO(6)$ 理论正是对纯 $\Spin(6)$ Yang--Mills 理论的 $\Z_4$ 电 1-形式对称性的 $\Z_2$ 子群作规范化而得到的。该反常的非平凡性与短正合列
\begin{equation}
  0\to\Z_2\to\Z_4\to\Z_2\to 0\,, \tag{4.199}
\end{equation}
不分裂这一事实有关。
\end{examplebox}

% [Tutor Note: p.95 的 “example (4.8)” 与例 4.16 中的 “statement (4.6)” 都是原文的引用笔误——括号是误加的：它们分别指例 4.8（Pure SO(6) Yang-Mills Theory）与陈述 4.6，原文他处引用时写作 “example 4.8” 或 “statement 4.6”。核对：p.95 所说“短正合列分裂则反常消失”，正是例 4.8 之前 “Vanishing of Bockstein: Split Short Exact Sequence” 一段（(4.114)--(4.118)，p.82）的结论。译文照原文保留括号。]

\subsection{自发对称性破缺与禁闭}

\subsubsection{一般概念}

正如普通的 0-形式整体对称性一样，高阶 p-形式整体对称性也可以表现出自发对称性破缺（SSB）。我们将看到，自发破缺与\textbf{禁闭}密切相关。

% [原文 p.96]

\para{0-形式对称性的 SSB。}

我们先把 0-形式对称性的自发破缺重新表述成一种可以推广到 p-形式对称性的形式。传统定义如下。

\begin{definitionbox}{定义 4.8：0-形式对称性的自发破缺}
若存在满足以下两个条件的局域算符 $\cO$，则称 0-形式对称群 $G^{(0)}$ 发生自发破缺：
\begin{enumerate}[leftmargin=1.4em,itemsep=1pt,topsep=1pt]
\item $\cO$ 在 $G^{(0)}$ 下带非平凡的荷；
\item $\cO$ 的 vev 非零
\begin{equation}
  \langle\cO\rangle\neq 0\,. \tag{4.200}
\end{equation}
\end{enumerate}
\end{definitionbox}

主要要求 (4.200) 可以等价地表述为 $\cO$ 的一个渐近两点关联函数非零：
\begin{equation}
  \lim_{|x-y|\to\infty}\langle\cO(x)\cO(y)\rangle=\langle\cO(x)\rangle\langle\cO(y)\rangle=\langle\cO\rangle\langle\cO\rangle\neq 0\,, \tag{4.201}
\end{equation}
其中我们用集团分解把大距离处的两点关联函数化成一乘积的单点关联函数。

\para{p-形式对称性的 SSB。}

把 (4.201) 推广，p-形式对称性的自发破缺定义如下。

\begin{definitionbox}{定义 4.9：p-形式对称性的自发破缺}
若存在一个 p 维算符 $\cO_p$，它在 $G^{(p)}$ 下带非平凡的荷，且具有渐近关联函数
\begin{equation}
  \lim_{R\to\infty}\langle\cO_p(S^p_R)\rangle\neq 0\,, \tag{4.202}
\end{equation}
其中 $\cO_p$ 沿 $S^p_R$ 插入、$S^p_R$ 是半径为 $R$ 的 p 维球面，则称 p-形式对称群 $G^{(p)}$ 发生自发破缺。
\end{definitionbox}

注意对 $p=0$ 这与 (4.201) 相符，因为 0 维球面由两个点组成。

\para{与禁闭的关系。}

满足 (4.202) 的这种算符 $\cO_p$ 在深 IR 中存活。另一方面，若某算符 $\cO'_p$ 满足
\begin{equation}
  \lim_{R\to\infty}\left\langle\cO'_p(S^p_R)\right\rangle=0\,, \tag{4.203}
\end{equation}
则算符 $\cO_p$ 在 IR 中不存在。

% [Tutor Note: p.96 末句 “then the operator $\cO_p$ does not exist in the IR” 中的 $\cO_p$ 应为 $\cO'_p$（原文笔误）：(4.203) 的主语是 $\cO'_p$，且 p.97 首句明确说“满足 (4.203) 的算符 $\cO'_p$ 禁闭”。判断依据：若按字面理解，$\cO'_p$ 的 vev 为零将推出另一算符 $\cO_p$ 不存在，与 (4.202) 的定义自相矛盾。译文照原文翻译，此注给出正确形式。]

% [原文 p.97]

因此，人们说满足 (4.203) 的算符 $\cO'_p$ 在 IR 中\textbf{禁闭}，而满足 (4.202) 的算符 $\cO_p$ 在 IR 中\textbf{去禁闭}。

于是我们得到如下陈述。

\begin{statementbox}{陈述 4.8：（去）禁闭与 SSB}
若一个在 p-形式对称性 $G^{(p)}$ 下带电的 p 维算符去禁闭，则 $G^{(p)}$ 自发破缺。

另一方面，若所有带电的 p 维算符都禁闭，则 $G^{(p)}$ 不发生自发破缺。
\end{statementbox}

\para{与周长律和面积律的关系。}

在讨论线算符的语境中，禁闭与去禁闭还有另一个著名的定义如下：

\begin{itemize}[leftmargin=1.4em,itemsep=4pt,topsep=1pt]
\item 若线算符 $L$ 的渐近 vev 表现为
\begin{equation}
  \lim_{R\to\infty}\langle L(S^1_R)\rangle\sim\exp\left(-\alpha R\right)\,, \tag{4.204}
\end{equation}
则称线 $L$ 呈现\textbf{周长律}。在这种情况下，$L$ 去禁闭。
\item 另一方面，若线算符 $L$ 的渐近 vev 表现为
\begin{equation}
  \lim_{R\to\infty}\langle L(S^1_R)\rangle\sim\exp\left(-\alpha R^2\right)\,, \tag{4.205}
\end{equation}
则称线 $L$ 呈现\textbf{面积律}。在这种情况下，$L$ 禁闭。
\end{itemize}

粗略看来，周长律与面积律似乎都满足禁闭的条件 (4.203)。然而关键在于：对于一个满足周长律 (4.204) 的线算符 $L$，我们可以沿该线加一个局域的抵消项来修改它，从而得到一个新线算符 $L'$
\begin{equation}
  L'(S^1_R):=L(S^1_R)\times\exp\left(\int_{S^1_R}\vol\right)\,, \tag{4.206}
\end{equation}
其中 $\vol$ 是 $S^1_R$ 上的体积形式。算符 $L'$ 现在表现出 (4.202)，从而去禁闭。于是人们说原来的线 $L$ 去禁闭——在心里保留上述理解即可。

另一方面，人们无法对一个满足面积律 (4.205) 的线算符 $L$ 加上类似的局域抵消项：

\begin{itemize}[leftmargin=1.4em,itemsep=4pt,topsep=1pt]
\item 若沿该线加一个局域抵消项，则所得的线算符 $L'$ 仍然禁闭。
\end{itemize}

% [Tutor Note: p.97 (4.206) 的解释：周长律 $\langle L\rangle\sim e^{-\alpha R}$ 只含 $R$ 的一次幂，而沿 $S^1_R$ 的积分 $\int\vol\propto R$，故可把指数从 $-\alpha R$ 调成常数（取合适的抵消项系数），使 $L'$ 满足 (4.202)。面积律 $e^{-\alpha R^2}$ 含 $R^2$，一维积分只能产生 $R$ 的线性项，无法抵消，见 p.98 的 (4.207) 必须用二维面 $\Sigma_2$（面积 $\sim R^2$），代价是 $L'$ 变成依附在面算符上的非真线算符。]

% [原文 p.98]

\begin{figure}[htbp]
\centering
\includegraphics[width=0.52\textwidth]{fig20.png}
\figtitle{20}{沿时空中一个矩形插入的线算符 $L$，其沿空间方向的长度为 $R$，沿时间方向的长度为 $T$。}
\end{figure}

\begin{itemize}[leftmargin=1.4em,itemsep=4pt,topsep=1pt]
\item 若想要一条去禁闭的线 $L'$，抵消项必须定义在一个面 $\Sigma_2$ 上，而 $\Sigma_2$ 的边界是线算符 $L$ 的位置，例如
\begin{equation}
  L'(S^1_R,\Sigma_2):=L(S^1_R)\times\exp\left(\int_{\Sigma_2}\vol\right)\,, \tag{4.207}
\end{equation}
其中 $\partial\Sigma_2=S^1_R$、$\vol$ 是 $\Sigma_2$ 上的体积形式。然而在这种情形下，所得的线算符 $L'$ 依赖于 $\Sigma_2$，因此它是依附在一个面算符上的非真线算符。
\end{itemize}

\para{探针粒子之间的势。}

若把线算符看作探针粒子，则渐近 vev 可以用这些探针粒子之间的势能来刻画。为此，考虑沿图 20 所示的矩形插入一个线算符 $L$。这个位形现在可以看作两个探针粒子相距 $R$、并保持时间间隔 $T$ 的振幅。于是 vev 可以计算为
\begin{equation}
  \langle L\rangle\sim\exp\left(-V(R)\,T\right)\,, \tag{4.208}
\end{equation}
其中 $V(R)$ 是探针粒子之间的势能。

在这种刻画下，周长律对应于
\begin{equation}
  V(R)\sim\text{常数}\,, \tag{4.209}
\end{equation}
而面积律对应于
\begin{equation}
  V(R)\sim R\,. \tag{4.210}
\end{equation}

因此，对于呈现面积律的探针粒子，把它们分开到无穷远需要无穷大的能量。换句话说，这些探针粒子保持禁闭。

% [Tutor Note: p.98 (4.208)--(4.210) 的来源：矩形环的周长 $\propto 2(R+T)$，代入 $\langle L\rangle=\exp(-V(R)T)$ 得 $V(R)=2\alpha+2\alpha R/T$，$T\to\infty$ 时趋于常数，即 (4.209)；而面积 $\propto RT$ 代入得 $V(R)=\alpha R$，即 (4.210)。注意 (4.204) 为书写简洁只写了 $R$，完整的周长律应同时含 $T$。]
